\documentclass[a4paper,11pt]{article}
\usepackage[margin=25mm]{geometry}
\usepackage{fontspec}
\usepackage{xeCJK}
\setCJKmainfont{Hiragino Mincho ProN}
\setCJKsansfont{Hiragino Sans}
\setCJKmonofont{Osaka-Mono}
\usepackage{hyperref}
\hypersetup{colorlinks=true,linkcolor=blue!60!black,urlcolor=blue!70!black}
\usepackage{booktabs}
\usepackage{longtable}
\usepackage{enumitem}
\usepackage{fancyhdr}
\usepackage{xcolor}
\usepackage{tcolorbox}
\usepackage{listings}
\usepackage{titlesec}
\usepackage{amssymb}

\definecolor{warn}{RGB}{180,60,20}
\definecolor{safe}{RGB}{20,120,60}
\definecolor{codebg}{RGB}{245,245,245}

\newtcolorbox{warningbox}{colback=red!5!white,colframe=red!60!black,title={\textbf{安全警告}},fonttitle=\sffamily}
\newtcolorbox{tipbox}{colback=blue!5!white,colframe=blue!40!black,title={\textbf{ヒント}},fonttitle=\sffamily}
\newtcolorbox{checklistbox}{colback=green!5!white,colframe=green!40!black,fonttitle=\sffamily}

\lstset{
  basicstyle=\ttfamily\small,
  backgroundcolor=\color{codebg},
  frame=single,
  framerule=0.5pt,
  rulecolor=\color{gray!50},
  breaklines=true,
  columns=fullflexible,
}

\pagestyle{fancy}
\fancyhf{}
\fancyhead[L]{\sffamily\small Open Excretion Care System --- 実践ビルドガイド}
\fancyhead[R]{\sffamily\small v0.1}
\fancyfoot[C]{\thepage}

\titleformat{\section}{\Large\sffamily\bfseries}{第\thesection 章}{1em}{}
\titleformat{\subsection}{\large\sffamily\bfseries}{\thesubsection}{1em}{}

\title{\vspace{-10mm}\sffamily\bfseries
Open Excretion Care System (OECS)\\[4pt]
\Large 実践ビルドガイド\\[2pt]
\large --- 排泄自動検知・洗浄システムを自作する ---}
\author{Based on Hu \& Chen et al., \textit{Healthcare} 2023, 11(3), 388}
\date{2026年4月}

\begin{document}
\maketitle
\thispagestyle{fancy}

\begin{warningbox}
本ガイドは教育・研究・個人利用を目的としています。医療機器としての認証は取得していません。実際の介護現場での使用前に、必ず医療従事者に相談し、各国の医療機器規制を確認してください。電気・水・人体接触を伴うため、安全対策を最優先にしてください。
\end{warningbox}

\vspace{3mm}
\tableofcontents
\newpage

%======================================================================
\section{プロジェクト概要}
%======================================================================

\subsection{何を作るのか}

本システムは、おむつ内に設置したセンサーで排泄を自動検知し、吸引$\to$温水洗浄$\to$温風乾燥のサイクルを自動実行する装置です。

\begin{itemize}[leftmargin=2em]
  \item \textbf{検知方式}: 温湿度＋アンモニアガスの多センサー融合（Dempster--Shafer証拠理論）
  \item \textbf{識別精度}: 約90\%（単純閾値方式の約70\%から大幅向上）
  \item \textbf{尿/便の識別}: ガスセンサーにより区別し、便の場合は洗浄時間を延長
  \item \textbf{推定部品コスト}: \$45--100 USD
\end{itemize}

\subsection{システム構成の全体像}

\begin{verbatim}
┌─────────────┐     ┌──────────────┐     ┌──────────────┐
│ センサー層   │ --> │  判定ロジック  │ --> │ アクチュエータ │
│ SHT30(温湿度)│     │  D-S証拠理論  │     │  吸引ポンプ   │
│ MQ-135(NH3) │     │  (or 1D-CNN) │     │  温水ポンプ   │
└─────────────┘     └──────────────┘     │  ヒーター+ファン│
       │                   │              └──────────────┘
       └───────────────────┴──────────────────────┘
                    STM32F103 MCU
\end{verbatim}

\subsection{本ガイドの構成}

\begin{enumerate}
  \item 部品の調達（BOM）
  \item 回路の組み立て
  \item レシーバーの3Dプリント
  \item ファームウェアの書き込み
  \item アルゴリズムの校正
  \item 組み立てと動作テスト
  \item AI拡張（オプション）
  \item トラブルシューティング
\end{enumerate}

\newpage
%======================================================================
\section{部品の調達}
%======================================================================

\subsection{必須部品リスト (BOM)}

\begin{longtable}{@{}clllr@{}}
\toprule
\textbf{\#} & \textbf{部品} & \textbf{型番/仕様} & \textbf{数量} & \textbf{目安価格} \\
\midrule
\endhead
1  & マイコンボード     & STM32F103C8T6 (Blue Pill)  & 1 & \$3  \\
2  & 温湿度センサー     & SHT30-DIS-B (I2C)         & 1 & \$5  \\
3  & ガスセンサー       & MQ-135モジュール            & 1 & \$3  \\
4  & 吸引ポンプ         & 12Vダイアフラムポンプ       & 1 & \$12 \\
5  & 温水ポンプ         & 12Vペリスタルティックポンプ  & 1 & \$10 \\
6  & ヒーター           & PTC セラミックヒーター 12V 50W & 1 & \$5  \\
7  & 送風ファン         & 5015ブロワーファン 12V       & 1 & \$3  \\
8  & モータードライバ   & L298N デュアルHブリッジ      & 1 & \$3  \\
9  & MOSFETモジュール   & IRF520                     & 2 & \$2  \\
10 & 電源アダプタ       & 12V 5A スイッチング         & 1 & \$8  \\
11 & 電圧レギュレータ   & AMS1117-3.3V               & 1 & \$1  \\
12 & 温度ヒューズ       & 73°C サーマルカットオフ      & 1 & \$1  \\
13 & 圧力リリーフバルブ & 調整式 $-20$\,kPa           & 1 & \$5  \\
14 & シリコンチューブ   & 食品グレード 内径6mm 2m      & 1 & \$4  \\
15 & TPUフィラメント     & Shore 60A 1.75mm 500g      & 1 & \$25 \\
16 & 配線キット         & 22AWG シリコン被覆          & 1 & \$5  \\
17 & コネクタ           & JST-XH 2.54mm              & 1 & \$3  \\
18 & 基板               & ユニバーサル基板 7$\times$9cm & 1 & \$2  \\
19 & 筐体               & PLA 3Dプリント              & 1 & \$3  \\
20 & 緊急停止ボタン     & NC プッシュボタン 赤 12mm   & 1 & \$2  \\
\midrule
   & & & \textbf{合計} & \textbf{\$\,\raisebox{0pt}{\~{}}103} \\
\bottomrule
\end{longtable}

\subsection{調達先の目安}

\begin{itemize}[leftmargin=2em]
  \item \textbf{AliExpress}: STM32 Blue Pill, MQ-135, IRF520, L298N, ポンプ類（安価だが納期2--4週間）
  \item \textbf{Amazon}: 電源アダプタ、シリコンチューブ、配線（即日配送可）
  \item \textbf{DigiKey/Mouser}: SHT30、温度ヒューズ（高精度部品、データシート確認用）
  \item \textbf{LCSC/JLCPCB}: PCBを本格製作する場合
\end{itemize}

\begin{tipbox}
最初のプロトタイプでは、SHT30の代わりにDHT22（\$2）でも始められます。精度は劣りますが、動作確認には十分です。
\end{tipbox}

\newpage
%======================================================================
\section{回路の組み立て}
%======================================================================

\subsection{ピンアサイン表}

\begin{longtable}{@{}llll@{}}
\toprule
\textbf{STM32ピン} & \textbf{機能} & \textbf{接続先} & \textbf{備考} \\
\midrule
\endhead
PA0  & ADC1\_CH0   & MQ-135 アナログ出力 & 0--3.3V \\
PA8  & GPIO出力    & L298N IN1（吸引ポンプ） & \\
PA9  & GPIO出力    & L298N IN3（温水ポンプ） & ※開発時はUART TX \\
PA10 & GPIO出力    & IRF520 \#1（ヒーター）  & ※開発時はUART RX \\
PA11 & GPIO出力    & IRF520 \#2（ファン）    & \\
PA12 & GPIO入力    & 緊急停止ボタン          & 外部プルダウン \\
PB6  & I2C1\_SCL   & SHT30 SCL & 4.7k$\Omega$プルアップ$\to$3.3V \\
PB7  & I2C1\_SDA   & SHT30 SDA & 4.7k$\Omega$プルアップ$\to$3.3V \\
PC13 & GPIO出力    & ステータスLED & オンボード（アクティブLow） \\
\bottomrule
\end{longtable}

\subsection{電源配線}

\begin{verbatim}
12V 5A ACアダプタ
├── L298N VCC (12V)
│   ├── OUT1/2 → 吸引ポンプ
│   └── OUT3/4 → 温水ポンプ
├── IRF520 #1 VCC (12V)
│   └── PTC ヒーター ── 73℃温度ヒューズ（直列）
├── IRF520 #2 VCC (12V)
│   └── ブロワーファン
└── AMS1117-3.3V
    ├── STM32 3.3V
    ├── SHT30 VCC
    └── (MQ-135ヒーターは5Vが必要 → L298Nの5V出力を利用)
\end{verbatim}

\subsection{組み立て手順}

\begin{enumerate}[leftmargin=2em]
  \item \textbf{電源系を最初に}: AMS1117-3.3Vレギュレータを基板に実装し、12V入力から3.3Vが正しく出力されることをテスターで確認。
  \item \textbf{MCUの動作確認}: Blue PillにUSB経由で給電し、PC13のLEDが点滅するテストファームを書き込む。
  \item \textbf{センサー接続}:
    \begin{itemize}
      \item SHT30: I2Cバス（PB6/PB7）に接続。4.7k$\Omega$プルアップ抵抗を忘れずに。
      \item MQ-135: アナログ出力をPA0に接続。\textbf{初回は24時間のバーンイン（通電エージング）が必要。}
    \end{itemize}
  \item \textbf{アクチュエータ接続}: L298NとIRF520を接続し、各モーター/ヒーター/ファンが個別にON/OFFできることを確認。
  \item \textbf{安全装置}: 緊急停止ボタンとサーマルヒューズを最後に配線。
\end{enumerate}

\begin{warningbox}
\textbf{必ず守ること}:
\begin{itemize}[leftmargin=1.5em]
  \item ヒーターには必ず73°Cサーマルヒューズを直列に入れる
  \item 吸引系には圧力リリーフバルブ（$-20$\,kPa）を取り付ける
  \item 温水の温度は38°C以下に制限（火傷防止）
  \item すべてのアクチュエータをOFFにしてから配線作業を行う
\end{itemize}
\end{warningbox}

\newpage
%======================================================================
\section{レシーバーの3Dプリント}
%======================================================================

\subsection{設計コンセプト}

レシーバーは患者の肌に直接触れる部品であり、以下の要件を満たす必要があります：

\begin{itemize}[leftmargin=2em]
  \item \textbf{素材}: TPU（Shore 60A）--- 柔軟で肌への負担が少ない
  \item \textbf{気密性}: リップシール構造で漏れを防止
  \item \textbf{ポート}: 吸引口（6mmバーブ）、給水口（6mmバーブ）、センサーポケット
  \item \textbf{サイズ展開}: S/M/Lのパラメトリック設計
\end{itemize}

\subsection{プリント設定}

\begin{longtable}{@{}ll@{}}
\toprule
\textbf{パラメータ} & \textbf{推奨値} \\
\midrule
レイヤー高さ & 0.2mm \\
インフィル & 20\% ジャイロイド \\
ウォール数 & 3 \\
プリント速度 & 25mm/s（TPUは遅めに） \\
リトラクション & ダイレクトドライブ推奨 \\
ベッド温度 & 50°C \\
ノズル温度 & 220--230°C \\
サポート & ツリーサポート \\
\bottomrule
\end{longtable}

\subsection{筐体（コントローラケース）}

\begin{itemize}[leftmargin=2em]
  \item \textbf{素材}: PLAまたはPETG
  \item \textbf{設計}: 換気スロット付き、メンテナンス用の脱着式蓋
  \item \textbf{プリント}: 0.2mm層高、30\%インフィル、60mm/s
  \item STLファイルは \texttt{hardware/3d-models/} に配置
\end{itemize}

\begin{tipbox}
TPU印刷に不慣れな場合、まずPLAでレシーバーのモックアップを印刷し、フィッティングを確認してからTPUで本番印刷すると材料の無駄を減らせます。
\end{tipbox}

\newpage
%======================================================================
\section{ファームウェアの書き込み}
%======================================================================

\subsection{開発環境のセットアップ}

\textbf{方法A: STM32CubeIDE（推奨・初心者向け）}

\begin{enumerate}[leftmargin=2em]
  \item ST公式サイトからSTM32CubeIDEをダウンロード・インストール
  \item 新規STM32プロジェクトを作成（ターゲット: STM32F103C8T6）
  \item CubeMXでピン設定をセクション3のピンアサイン通りに設定
  \item 生成されたプロジェクトに本リポジトリの \texttt{firmware/Core/} と \texttt{algorithm/c/} のファイルをコピー
\end{enumerate}

\textbf{方法B: コマンドライン（上級者向け）}

\begin{lstlisting}[language=bash]
# ARM GCCツールチェーンのインストール
# macOS:
brew install arm-none-eabi-gcc
# Ubuntu:
sudo apt install gcc-arm-none-eabi

# ビルド
cd firmware
make

# 書き込み (ST-Linkアダプタ使用)
make flash
\end{lstlisting}

\subsection{ST-Linkでの書き込み}

\begin{enumerate}[leftmargin=2em]
  \item ST-Link V2（クローン可、\$3程度）をBlue Pillに接続:
    \begin{itemize}
      \item SWDIO $\to$ PA13
      \item SWCLK $\to$ PA14
      \item GND $\to$ GND
      \item 3.3V $\to$ 3.3V（USBから給電しない場合のみ）
    \end{itemize}
  \item Blue PillのBOOT0ジャンパを0（通常動作）に設定
  \item \texttt{st-flash write build/oecs\_firmware.bin 0x8000000} で書き込み
\end{enumerate}

\subsection{デバッグ出力の確認}

UART1（115200bps）をUSB-シリアルアダプタ経由でPCに接続すると、以下のようなログが表示されます：

\begin{lstlisting}
=== OECS Firmware v0.1 ===
Open Excretion Care System initialized.
[DETECT] URINE (T=35.8 H=88.0 NH3=18.0)
[CYCLE] Starting cleaning cycle
[PHASE] Suction complete -> Washing
[PHASE] Washing complete -> Drying
[PHASE] Drying complete -> Cycle finished
[CYCLE] Returning to monitoring
\end{lstlisting}

\newpage
%======================================================================
\section{アルゴリズムの校正}
%======================================================================

\subsection{D-S証拠理論の概要}

本システムは、単純な閾値判定ではなく、\textbf{Dempster--Shafer（D-S）証拠理論}による多センサー融合を使用しています。

\textbf{処理フロー}:

\begin{verbatim}
温度データ → temp_to_bpa() → m_temp ─┐
                                       ├→ ds_combine() → m_12 ─┐
湿度データ → hum_to_bpa()  → m_hum  ─┘                         │
                                                                 ├→ ds_combine() → 最終判定
NH3データ  → nh3_to_bpa()  → m_nh3  ───────────────────────────┘
\end{verbatim}

各センサーから「正常/尿/便である確率」を算出し、D-S合成則で統合します。

\subsection{基本確率割り当て（BPA）テーブル}

\textbf{温度（基準: 33°C）}:
\begin{center}
\begin{tabular}{lcccc}
\toprule
\textbf{変化量} & \textbf{正常} & \textbf{尿} & \textbf{便} & \textbf{不確実} \\
\midrule
$< 0.5$°C  & 0.70 & 0.05 & 0.05 & 0.20 \\
$0.5$--$2.0$°C & 0.30 & 0.20 & 0.20 & 0.30 \\
$2.0$--$4.0$°C & 0.05 & 0.50 & 0.20 & 0.25 \\
$> 4.0$°C  & 0.02 & 0.55 & 0.25 & 0.18 \\
\bottomrule
\end{tabular}
\end{center}

\textbf{湿度}:
\begin{center}
\begin{tabular}{lcccc}
\toprule
\textbf{湿度} & \textbf{正常} & \textbf{尿} & \textbf{便} & \textbf{不確実} \\
\midrule
$< 60\%$    & 0.75 & 0.05 & 0.05 & 0.15 \\
$60$--$75\%$  & 0.30 & 0.25 & 0.15 & 0.30 \\
$75$--$85\%$  & 0.05 & 0.45 & 0.25 & 0.25 \\
$> 85\%$    & 0.02 & 0.60 & 0.20 & 0.18 \\
\bottomrule
\end{tabular}
\end{center}

\textbf{NH3ガス}:
\begin{center}
\begin{tabular}{lcccc}
\toprule
\textbf{NH3濃度} & \textbf{正常} & \textbf{尿} & \textbf{便} & \textbf{不確実} \\
\midrule
$< 5$\,ppm    & 0.80 & 0.05 & 0.02 & 0.13 \\
$5$--$15$\,ppm  & 0.10 & 0.50 & 0.15 & 0.25 \\
$15$--$30$\,ppm & 0.03 & 0.50 & 0.25 & 0.22 \\
$> 30$\,ppm   & 0.02 & 0.10 & 0.70 & 0.18 \\
\bottomrule
\end{tabular}
\end{center}

\subsection{Pythonシミュレータでの校正}

実際のセンサーを接続する前に、Python実装で動作を確認できます。

\begin{lstlisting}[language=bash]
cd algorithm/python
pip install numpy

# デモ実行（6パターンのテストケース）
python ds_evidence.py --demo

# カスタム値でテスト
python ds_evidence.py --temp 35.2 --humidity 85.0 --nh3 18.0
\end{lstlisting}

\textbf{デモ出力例}:

\begin{lstlisting}
--- Urine event ---
  Input:  T=35.8 C  RH=88.0%  NH3=18.0ppm
  Belief: Normal=0.010  Urine=0.828  Feces=0.142  Uncertain=0.020
  Result: URINE
  Action: Urine detected: suction -> rinse -> dry

--- Feces event ---
  Input:  T=35.0 C  RH=82.0%  NH3=42.0ppm
  Belief: Normal=0.020  Urine=0.399  Feces=0.548  Uncertain=0.033
  Result: FECES
  Action: Feces detected: extended suction -> thorough rinse -> dry
\end{lstlisting}

\subsection{校正の手順}

実環境でのBPA値チューニング手順：

\begin{enumerate}[leftmargin=2em]
  \item \textbf{ベースライン測定}: センサーをおむつ内に設置し、30分間のNormal状態を記録。温度・湿度・NH3の平均値と分散を取得。
  \item \textbf{尿イベント記録}: 尿が発生した際のセンサー変化パターンを10件以上記録。
  \item \textbf{便イベント記録}: 便が発生した際のパターンを同様に記録。
  \item \textbf{BPA調整}: 記録データに基づき、\texttt{ds\_evidence.py}内の \texttt{temperature\_to\_bpa()}, \texttt{humidity\_to\_bpa()}, \texttt{nh3\_to\_bpa()} 各関数の閾値と確率値を調整。
  \item \textbf{C実装への反映}: Python側で確定した値を \texttt{algorithm/c/ds\_evidence.c} のQ16固定小数点定数に変換。
\end{enumerate}

\begin{tipbox}
Q16固定小数点への変換: 小数値 $\times$ 65536 で整数化。例: $0.70 \times 65536 = 45875$
\end{tipbox}

\newpage
%======================================================================
\section{組み立てと動作テスト}
%======================================================================

\subsection{組み立て順序}

\begin{enumerate}[leftmargin=2em]
  \item \textbf{Phase 1: 電子系の確認}（机上テスト）
    \begin{itemize}
      \item ファームウェアを書き込んだBlue Pillに、SHT30とMQ-135を接続
      \item UART出力でセンサー値が正常に読めることを確認
      \item アクチュエータ（ポンプ・ヒーター・ファン）を個別にON/OFFテスト
    \end{itemize}
  \item \textbf{Phase 2: 水系の確認}
    \begin{itemize}
      \item シリコンチューブでポンプを接続し、水の吸引・送水を確認
      \item 温水温度が38°C以下であることを温度計で実測
      \item 圧力リリーフバルブの動作を確認
    \end{itemize}
  \item \textbf{Phase 3: レシーバー統合}
    \begin{itemize}
      \item 3DプリントしたTPUレシーバーにチューブとセンサーを取り付け
      \item 気密性テスト（水を入れて漏れ確認）
    \end{itemize}
  \item \textbf{Phase 4: 統合テスト}
    \begin{itemize}
      \item レシーバーに温水を注入し（尿シミュレート）、自動検知$\to$吸引$\to$洗浄$\to$乾燥の全サイクルが動くことを確認
      \item 緊急停止ボタンですべてのアクチュエータが即座にOFFになることを確認
    \end{itemize}
\end{enumerate}

\subsection{テストチェックリスト}

\begin{checklistbox}
\begin{itemize}[leftmargin=1.5em, label=$\square$]
  \item SHT30から温度・湿度が読める（UART出力で確認）
  \item MQ-135から安定したADC値が読める（24時間バーンイン後）
  \item 吸引ポンプがON/OFFできる
  \item 温水ポンプがON/OFFできる
  \item ヒーターがON/OFFできる（温度ヒューズ通電確認済み）
  \item ファンがON/OFFできる
  \item 温水を注入した際に「URINE」として検知される
  \item NH3源を近づけた際に「FECES」として検知される
  \item 正常時に誤検知しない（30分間の安定稼働）
  \item 緊急停止ボタンで全アクチュエータが即時停止する
  \item 洗浄水温が38°C以下である
  \item 吸引圧力がリリーフバルブで$-20$\,kPa以下に制限されている
  \item 全サイクル（吸引$\to$洗浄$\to$乾燥）が自動で完了する
\end{itemize}
\end{checklistbox}

\newpage
%======================================================================
\section{AI拡張（オプション）}
%======================================================================

\subsection{D-S理論 vs 1D-CNN}

\begin{center}
\begin{tabular}{lll}
\toprule
& \textbf{D-S証拠理論} & \textbf{1D-CNN} \\
\midrule
精度         & $\sim$90\%     & $>$95\%（十分なデータがあれば） \\
学習データ   & 不要           & 必要（100件以上推奨） \\
計算コスト   & 極めて低い     & 中程度（STM32でも動作可） \\
解釈性       & 高い           & 低い（ブラックボックス） \\
時系列認識   & なし           & あり（パターン認識可能） \\
予測機能     & なし           & 可能（排泄パターン学習） \\
\bottomrule
\end{tabular}
\end{center}

\subsection{1D-CNNモデルの概要}

\begin{verbatim}
入力: (batch, 3ch, 60 samples) -- 温度/湿度/NH3の60秒分の時系列
  ↓
Conv1D(3→16, k=5) → BatchNorm → ReLU → MaxPool(2)
  ↓
Conv1D(16→32, k=3) → BatchNorm → ReLU → MaxPool(2)
  ↓
Conv1D(32→32, k=3) → BatchNorm → ReLU → GlobalAvgPool
  ↓
FC(32→16) → ReLU → Dropout(0.3) → FC(16→3)
  ↓
出力: [P(正常), P(尿), P(便)]

総パラメータ数: 約3,500（STM32のSRAMに余裕で収まる）
\end{verbatim}

\subsection{トレーニング手順}

\begin{lstlisting}[language=bash]
cd ai/training
pip install torch numpy onnx

# 合成データで学習（初回テスト用）
python train_1dcnn.py --synthetic --epochs 100

# 学習済みモデルをONNXエクスポート
# → oecs_model.onnx が生成される

# STM32への展開:
# 1. ST X-CUBE-AI でONNXをCコードに変換
# 2. または TensorFlow Lite Micro で量子化→展開
\end{lstlisting}

\subsection{実データ収集のコツ}

\begin{enumerate}[leftmargin=2em]
  \item D-S理論ベースで最初のプロトタイプを稼働させる
  \item UART経由でセンサーデータ（1秒間隔）をPCにログ出力
  \item 介護者が排泄イベント時にタイムスタンプを記録（ラベリング）
  \item 100件以上のラベル付きデータが集まったらCNNの学習を試みる
\end{enumerate}

\newpage
%======================================================================
\section{トラブルシューティング}
%======================================================================

\begin{longtable}{@{}p{4.5cm}p{4.5cm}p{5cm}@{}}
\toprule
\textbf{症状} & \textbf{原因の可能性} & \textbf{対処法} \\
\midrule
\endhead
センサー値が読めない & I2C配線ミス、プルアップ抵抗なし & PB6/PB7の配線確認、4.7k$\Omega$追加 \\
\addlinespace
MQ-135の値が不安定 & バーンイン不足 & 24時間以上通電してエージング \\
\addlinespace
常に「NORMAL」判定 & センサーがレシーバー内にない、閾値が高すぎる & センサー位置確認、BPA閾値を下げる \\
\addlinespace
誤検知が多い & 汗・体温による誤反応 & デバウンスカウントを増やす（3→5）、閾値を上げる \\
\addlinespace
ポンプが動かない & L298Nのイネーブルピン未接続 & ENAとENBをHIGH（ジャンパ接続）に \\
\addlinespace
温水が熱すぎる & ヒーター制御の問題 & サーマルヒューズの動作確認、PWMデューティ調整 \\
\addlinespace
吸引力が弱い & チューブの折れ、リーク & チューブの取り回し確認、接続部のシール \\
\addlinespace
ファームが書き込めない & BOOT0設定、ST-Link接続 & BOOT0=0確認、SWDIO/SWCLKの結線確認 \\
\addlinespace
緊急停止が効かない & \textcolor{warn}{\textbf{直ちに電源を切る}} & GPIO入力の極性確認、ハードワイヤード回路の追加を検討 \\
\bottomrule
\end{longtable}

\newpage
%======================================================================
\section{安全設計の補足}
%======================================================================

\begin{warningbox}
本システムは\textbf{水・電気・人体}が関わります。以下の安全対策は\textbf{省略不可}です。
\end{warningbox}

\subsection{多重安全機構}

\begin{enumerate}[leftmargin=2em]
  \item \textbf{ソフトウェア制限}: ファームウェア内で温水38°C上限、吸引圧$-20$\,kPa上限
  \item \textbf{ハードウェア制限}: サーマルヒューズ（73°C）、圧力リリーフバルブ
  \item \textbf{物理的遮断}: 緊急停止ボタン（電源ライン直接遮断を推奨）
  \item \textbf{電気的隔離}: 患者接触部分と高電圧部分はオプトカプラで絶縁
\end{enumerate}

\subsection{推奨される追加安全策}

\begin{itemize}[leftmargin=2em]
  \item 漏電ブレーカーの設置（感電防止）
  \item 水系統にIPX4以上の防水処理
  \item 定期的な気密性テスト（週1回）
  \item センサーの定期校正（月1回）
\end{itemize}

\vspace{10mm}

%======================================================================
\section{付録: リソースとリンク}
%======================================================================

\begin{itemize}[leftmargin=2em]
  \item \textbf{原論文}: Hu, B.; Chen, Z.; et al. ``Research of System Design and Automatic Detection Method for Excretion Nursing Equipment.'' \textit{Healthcare} 2023, 11(3), 388. \\ \url{https://doi.org/10.3390/healthcare11030388}
  \item \textbf{GitHubリポジトリ}: \texttt{excretion-care-system/} --- 本ガイドに対応するすべてのソースコード
  \item \textbf{STM32 HALドキュメント}: ST公式 STM32CubeF1パッケージ
  \item \textbf{SHT30データシート}: Sensirion公式
  \item \textbf{MQ-135データシート}: Winsen公式
  \item \textbf{D-S証拠理論}: Shafer, G. \textit{A Mathematical Theory of Evidence}. Princeton UP, 1976.
\end{itemize}

\vfill
\begin{center}
\rule{0.5\textwidth}{0.4pt}\\[6pt]
{\small 本ガイドはオープンソースです。改善提案やフィードバックをお待ちしています。}\\
{\small CERN-OHL-S-2.0 (Hardware) / GPLv3 (Software)}
\end{center}

\end{document}
